%%
%% FedCSL-Spilter -- ACM SIGMOD-style conference article.
%% Derived from the ACM Primary Article Template (sigconf.tex), renamed per the
%% template notice ("Any modified versions of this file must be renamed").
%% Compile with:  tectonic -X compile fedcsl_spilter_acm.tex
%%
\documentclass[sigconf]{acmart}

\AtBeginDocument{%
  \providecommand\BibTeX{{%
    Bib\TeX}}}

%% Rights information: replace with the values from the rights confirmation
%% e-mail once the paper is accepted.
\setcopyright{acmlicensed}
\copyrightyear{2026}
\acmYear{2026}
\acmDOI{XXXXXXX.XXXXXXX}
\acmConference[Submitted]{ACM Conference (under review)}{2026}{Harbin, China}
\acmISBN{978-1-4503-XXXX-X/2026/06}

\usepackage{booktabs}
\usepackage{amsmath}
\usepackage{amssymb}
\usepackage{graphicx}
\usepackage{multirow}
\usepackage{array}
\usepackage{algorithm}
\usepackage{algpseudocode}
\usepackage{enumitem}

\begin{document}

\title[FedCSL-Spilter]{FedCSL-Spilter: Federated Shapelet Representation Learning
with Adaptive Scale Stitching and Memory-Aware Planning for Non-IID
Multivariate Time Series}

\author{Xiongfei Li}
\affiliation{%
  \institution{Harbin Institute of Technology}
  \city{Harbin}
  \country{China}}
\email{author@example.com}   % TODO: replace with the real address before submission

\renewcommand{\shortauthors}{Li}

\begin{abstract}
Multivariate time series (MTS) are ubiquitous in human activity recognition,
industrial monitoring, and medical diagnosis, yet the data are naturally
distributed across terminals and institutions, label-scarce, statistically
heterogeneous, and privacy-sensitive. Federated learning (FL) keeps raw data
local, but under non-IID partitions it suffers from client drift, cross-client
representation inconsistency, and excessive communication, computation, and GPU
memory cost. Existing federated time-series methods focus mainly on supervised
classification, parameter averaging, or client scheduling; they offer limited
support for interpretable multi-scale shapelet submodel construction and for
memory-aware execution under heterogeneous devices.

We propose \textbf{FedCSL-Spilter}, a federated shapelet learning framework
with three cooperating levels of optimization. At the \emph{representation}
level, local instance contrastive learning is augmented with local--global joint
contrastive learning and joint knowledge distillation (KD), so that client models
retain local adaptability while staying aligned with a shared global semantic
space. At the \emph{scale-composition} level, \textbf{Spilter} lets every client
rank $R$ candidate shapelet scales by STFT/ACF period-aware scores computed on
local data, select the top-$m$ scales, stitch their shapelet features along the
representation dimension into a local submodel, upload only the parameters of
trained scales, and lets the server aggregate in a scale-wise manner. At the
\emph{execution} level, \textbf{MAST} exploits the structural properties of the
shapelet encoder---activations dominated by independent per-scale
submodules---to cast activation-checkpoint planning as an exact binary linear
program, trading recomputation for memory without recursive dependencies.

We further formalize local scale selection as a covering-constrained 0--1
knapsack that a Lagrange multiplier decomposes across clients, so that global
scale coverage is preserved under per-device memory budgets. Experiments on five
MTS benchmarks---HAR, FaceDetection, LSST, Epilepsy, and SleepEDF---under
Dirichlet concentrations $\alpha\in\{0.05,\dots,10.0\}$ show that FedCSL-Spilter
matches or exceeds full-scale FedCSL and federated self-supervised baselines in
downstream SVM accuracy while cutting client peak memory by up to 71.1\%,
reducing the slowest-client epoch time by $1.48\times$, and, in the MAST
evaluation on four large UEA datasets, reducing training time by
1.43\%--7.65\% over state-of-the-art checkpoint planners at the same memory
budget.
\end{abstract}

\begin{CCSXML}
<ccs2012>
   <concept>
       <concept_id>10010147.10010257.10010293.10010294</concept_id>
       <concept_desc>Computing methodologies~Neural networks</concept_desc>
       <concept_significance>500</concept_significance>
   </concept>
   <concept>
       <concept_id>10010147.10010257.10010282.10011304</concept_id>
       <concept_desc>Computing methodologies~Learning paradigms</concept_desc>
       <concept_significance>500</concept_significance>
   </concept>
   <concept>
       <concept_id>10002951.10003227</concept_id>
       <concept_desc>Information systems~Data mining</concept_desc>
       <concept_significance>300</concept_significance>
   </concept>
   <concept>
       <concept_id>10011007.10011006.10011041</concept_id>
       <concept_desc>Software and its engineering~Distributed systems organizing principles</concept_desc>
       <concept_significance>300</concept_significance>
   </concept>
   <concept>
       <concept_id>10010147.10010341.10010349</concept_id>
       <concept_desc>Computing methodologies~Distributed artificial intelligence</concept_desc>
       <concept_significance>300</concept_significance>
   </concept>
</ccs2012>
\end{CCSXML}

\ccsdesc[500]{Computing methodologies~Neural networks}
\ccsdesc[500]{Computing methodologies~Learning paradigms}
\ccsdesc[300]{Information systems~Data mining}
\ccsdesc[300]{Software and its engineering~Distributed systems organizing principles}
\ccsdesc[300]{Computing methodologies~Distributed artificial intelligence}

\keywords{federated learning; multivariate time series; shapelet;
multi-scale stitching; adaptive scale selection; contrastive learning;
knowledge distillation; activation checkpointing}

\maketitle

%%--------------------------------------------------------------------
\section{Introduction}
%%--------------------------------------------------------------------

Time series analysis is a core topic in artificial intelligence and data mining,
with applications in smart manufacturing, EEG diagnosis, human activity
recognition, financial forecasting, and industrial alerting. Unlike static
tabular data, time series encode temporal dependencies, periodic structure, and
cross-scale patterns in addition to multivariate observations. Learning
discriminative, transferable, and interpretable temporal representations
therefore remains a central challenge.

Unsupervised representation learning has advanced rapidly on this problem.
Contrastive methods construct positive--negative relations across augmented
views and capture intrinsic consistency in feature space, improving downstream
classification, clustering, and anomaly detection~\cite{ts2vec,tstcc,tfc}.
Later work extends self-supervision through masked modelling~\cite{simmtm},
time--frequency consistency~\cite{tfc}, soft contrastive
relations~\cite{softclt}, and revisited positive/negative construction
strategies~\cite{timesurl}. TimeCSL/CSL further shows that
\emph{shapelets}---interpretable subsequence patterns---can be coupled with
contrastive learning to form a general unsupervised shapelet representation
framework~\cite{timescsl}, which is attractive for interpretability and
transferability~\cite{interpgn,unishape}. Most of these methods, however, use
CNN/Transformer backbones~\cite{fcnbaseline,timesnet} that were not designed for
variable-length, strongly periodic series.

In practice, large volumes of time-series data cannot be centralized: medical,
physiological, sensing, and edge data live on different clients and are
privacy-sensitive. Federated learning exchanges parameters or representations
instead of raw data~\cite{fedavg}, and proximal or contrastive
regularization~\cite{fedprox,moon} mitigates client drift. Federated
representation learning for multivariate time series---especially interpretable
shapelet learning---remains underexplored~\cite{fedst,fedtsc}.

Directly porting centralized TimeCSL/CSL to FL is non-trivial for three reasons.
First, client data are strongly non-IID across users, devices, and environments;
naive local training plus parameter averaging overfits local modes and drifts the
representation space. Second, time series exhibit rich multi-scale structure, and
treating all scales equally reduces the specificity of shapelet learning. Third,
``model heterogeneity'' here manifests as \emph{scale heterogeneity}: some scales
carry more parameters and heavier sliding-window matching than others, and
uniform or random scale usage wastes computation, communication, and GPU memory.

Prior heterogeneous FL methods adapt to device resources with width/depth
submodels (HeteroFL~\cite{heterofl}, FjORD~\cite{fjord}, DepthFL~\cite{depthfl},
ScaleFL~\cite{scalefl}), weak-device participation
(InclusiveFL~\cite{inclusivefl}, AdaptiveFL~\cite{adaptivefl}), or flexible
pruning with self-KD (FlexFL~\cite{flexfl}). These works show that submodel
structure is a first-order factor in FL efficiency, but they target generic
CNN/Transformer backbones and do not model \emph{scale cost} and \emph{periodic
utility} in multi-scale shapelet encoders. On the systems side, activation
checkpointing~\cite{chen2016checkpoint} reduces peak memory by recomputation, and
planners such as Checkmate~\cite{checkmate}, Mimose~\cite{mimose}, and
MONet~\cite{monet} search for good recomputation plans; generic planners either
lack global optimality (heuristics) or must approximate nonlinear dependency
constraints.

We propose \textbf{FedCSL-Spilter}, which couples federated shapelet
representation learning with client-autonomous \emph{scale stitching} and
memory-aware execution. The contributions are:

\begin{enumerate}
  \item A unified framework that extends interpretable multi-scale shapelet
        representation learning to non-IID FL while preserving TimeCSL/CSL's
        decoupled ``representation + analyzer'' paradigm.
  \item Local--global joint CL/KD and scale-level CL/KD that softly align client
        features to a global semantic space and suppress scale inconsistency
        under non-IID data.
  \item STFT/ACF period-aware weighting that adapts multi-scale losses to the
        dominant periodic structure of each client and supplies the utility
        signal for scale selection.
  \item \textbf{Spilter}: a purely local, distribution-aware top-$m$ scale
        selection and feature-stitching mechanism with scale-sharded
        communication and scale-wise aggregation, compared against random
        top-$m$, fixed scales, and full-scale FedCSL.
  \item A covering-constrained 0--1 knapsack formulation whose Lagrangian
        decomposition lets clients select scales independently while the server
        iteratively restores global scale coverage.
  \item \textbf{MAST}: an exact binary linear program for activation-checkpoint
        planning that exploits submodule independence in the shapelet encoder,
        used as a per-device memory--time regulator in synchronous FL.
  \item Experiments on five MTS benchmarks and four large UEA datasets covering
        accuracy, ablations, non-IID strength, convergence, interpretability,
        and system efficiency.
\end{enumerate}

%%--------------------------------------------------------------------
\section{Related Work}
%%--------------------------------------------------------------------

\subsection{Unsupervised Time-Series Representation Learning}

Unsupervised time-series representation learning seeks general temporal
embeddings without labels. Early reconstruction, forecasting, and autoencoder
approaches emphasize local reconstruction error but do not necessarily yield
discriminative downstream features~\cite{simmtm}. Contrastive learning maximizes
agreement across augmented views while separating negatives: TS2Vec performs
hierarchical contextual contrast~\cite{ts2vec}; TS-TCC combines temporal and
contextual contrast~\cite{tstcc}; TF-C enforces time--frequency
consistency~\cite{tfc}; SoftCLT mitigates over-repulsion of neighbouring
samples~\cite{softclt}; TimesURL revisits positive--negative construction and
SSCL losses~\cite{timesurl}; SimMTM extends self-supervision via masked
modelling and manifold neighbour reconstruction~\cite{simmtm}. A complementary
line builds models around interpretable subsequence patterns such as
shapelets~\cite{interpgn,unishape}; TimeCSL/CSL couples shapelets with
contrastive learning across scales and similarity
measures~\cite{timescsl}, enabling visualization and task transfer beyond
black-box embeddings.

\subsection{Federated Time-Series Learning and Knowledge Distillation}

Federated time-series learning must train collaboratively without sharing raw
data~\cite{fedavg}. Most prior work targets supervised classification through
parameter averaging, personalization, distillation, or transfer~\cite{fedkdsurvey}.
FedST proposes a secure federated shapelet
transformation~\cite{fedst}; FedTSC embeds interpretable TSC models in
FL~\cite{fedtsc}; EFDLS applies federated distillation to multi-task time-series
classification with feature-level teacher--student alignment~\cite{efdls}.
These results indicate that parameter averaging alone is insufficient under
strong temporal heterogeneity, and that feature-level distillation and semantic
alignment matter~\cite{moon,fedx}.

Cross-domain federated unsupervised representation learning is also active.
FedCA aligns cross-client representations for federated unsupervised
representation learning~\cite{fedca}; FedMoCo adapts MoCo to decentralized
medical imaging~\cite{fedmoco}; FedX uses bidirectional KD and local--global
contrast~\cite{fedx}; FedU2 identifies representation collapse and
representation-space inconsistency in federated unsupervised learning and gives
a unified aggregation scheme~\cite{fedu2}. On the supervised side, MOON adds
model-level contrast~\cite{moon}, FedProc and FedPCL construct contrastive
signals from class prototypes~\cite{fedproc,fedpcl}, and CDL suppresses
heterogeneity across vehicles~\cite{fedcdl}. FLESD performs similarity
distillation on the server side~\cite{flesd}; FLPD and FedRCIL offer
communication-efficient distillation variants~\cite{flpd,fedrcil}, while
FedCOG~\cite{fedcog}, MrTF~\cite{mrtf}, FedIOCD~\cite{fediocd}, and
FedNoRo~\cite{fednoro} broaden the boundary of federated distillation.

Most of these methods assume a supervised objective or visual unsupervised
learning; few study unsupervised shapelet representation directly. We target
general temporal representations in FL---learn an interpretable embedding space
first, then solve downstream tasks---which inherits the strengths of
TimeCSL~\cite{timescsl} and suits label-scarce settings~\cite{fedsmall}.

\subsection{Heterogeneous Federated Models and Scale-Level Composition}

Real FL systems cannot assume uniform client capacity: bandwidth, device
heterogeneity, and client scale constrain what each node can train. Oort and
FedCS show that statistical utility, device speed, bandwidth, and training
latency jointly determine time-to-accuracy~\cite{oort,fedcsadaptive}. We adopt
this systems view but do not study client sampling; instead we ask how, for a
given participant set, each node composes a scale submodel that fits its local
data.

Heterogeneous FL derives multiple submodels from one global model: HeteroFL
shrinks width~\cite{heterofl}; DepthFL varies depth~\cite{depthfl}; FjORD uses
ordered dropout to obtain nested submodels~\cite{fjord}; ScaleFL adapts width and
depth jointly~\cite{scalefl}; InclusiveFL and AdaptiveFL co-train large and small
models with distillation~\cite{inclusivefl,adaptivefl}; FlexFL estimates neuron
redundancy and prunes adaptively with self-KD~\cite{flexfl}. Efficiency depends
not only on participation but on whether submodel granularity matches device
resources. For FedCSL the natural structural unit is the \emph{shapelet scale}
rather than channels or layers, and clients with different periodic/spectral
structure naturally need different scales. We therefore instantiate heterogeneity
as client-autonomous \emph{scale stitching}, where device memory/compute
differences enter through per-client budgets $m_k$ rather than through the scale
ranking itself (Sec.~\ref{sec:resource-hetero}).

\subsection{Memory-Efficient Training and Checkpoint Planning}

Training deep networks stores intermediate activations for backpropagation, and
these activations---not parameters or optimizer states---dominate peak GPU memory
for long sequences and large batches. Activation checkpointing (AC) discards
selected activations during the forward pass and recomputes them during the
backward pass~\cite{chen2016checkpoint}, trading computation for memory. Sublinear
scheduling~\cite{chen2016checkpoint}, Mimose~\cite{mimose},
Checkmate~\cite{checkmate}, and MONet~\cite{monet} aim at good checkpoint plans; exactly solving the planning
problem is attractive because it removes tuning and guarantees budget
feasibility, but generic MILP formulations must encode nonlinear layer
dependencies and can be slow to solve for large networks. Our MAST planner
(Sec.~\ref{sec:mast}) restricts the formulation to the shapelet encoder's
independent submodules, which makes the program linear and the solution exact.

%%--------------------------------------------------------------------
\section{Problem Formulation}
%%--------------------------------------------------------------------

\subsection{Federated Setting}

Consider a central server and $K$ clients. Client $k$ holds local multivariate
time-series data
\begin{equation}
D_k = \{x_{k,i}\}_{i=1}^{n_k}, \qquad x_{k,i}\in\mathbb{R}^{C\times T},
\end{equation}
where $C$ is the number of channels, $T$ the sequence length, and $n_k$ the local
sample size; $N=\sum_{k=1}^{K} n_k$. Clients are typically non-IID,
\begin{equation}
P_k(x)\neq P_j(x),\quad k\neq j,
\end{equation}
reflecting differences in label distribution, temporal scale, periodicity,
amplitude, and local noise. We later instantiate heterogeneity with Dirichlet
partitions of concentration $\alpha$, where smaller $\alpha$ means stronger
heterogeneity.

\subsection{Objective}

We do not train an end-to-end classifier. Instead we learn a federated shapelet
encoder $f(\cdot;w)$ that maps a raw series to an interpretable representation,
\begin{equation}
z = f(x;w)\in\mathbb{R}^{d},
\end{equation}
and evaluate it with a downstream classifier (SVM), following TimeCSL's decoupled
``representation + task analyzer'' paradigm~\cite{timescsl}. The central
question is how to learn $w$ collaboratively when clients are non-IID and their
devices cannot all afford the full multi-scale model.

\subsection{Overall Pipeline}

Figure~\ref{fig:framework} summarizes FedCSL-Spilter. The server initializes the
global shapelet encoder. Each client precomputes period-aware scale utility
scores on its local data, selects the top-$m$ scales, and stitches their features
into a local submodel. Clients pull only the parameters of selected scales,
perform shapelet contrastive learning with local--global joint distillation and
scale-level distillation, and upload updates for trained scales. The server
aggregates scale-sharded updates with sample-size weighting, leaving untouched
scales unchanged. Period-aware weights from STFT/ACF statistics modulate the
multi-scale losses and provide the utility signal for scale selection. The loop
repeats until convergence, after which representations are evaluated downstream.

\begin{figure*}[t]
  \centering
  \includegraphics[width=0.92\textwidth]{figures/framework.png}
  \caption{Overall framework of FedCSL-Spilter. \textbf{Server (top):} a
  scale-wise aggregator fuses client updates per scale into the global shapelet
  encoder $w_g$; downstream evaluation trains an SVM on extracted
  representations. \textbf{Client $k$ (bottom):} multivariate series are
  augmented into two views $x_q,x_k$; the client selects $m$ scales and stitches
  shapelet features; an alignment node aggregates local and global features for
  $\mathcal{L}_{\text{local}}$, $\mathcal{L}_{\text{joint-CL}}$, and
  $\mathcal{L}_{\text{joint-KD}}$; the period-aware weighting module
  (bottom-right) produces scale weights $\boldsymbol{\pi}$ from an STFT branch
  (short--medium scales) and an ACF branch (medium--long scales), used both for
  multi-scale losses and for local scale selection.}
  \label{fig:framework}
\end{figure*}

%%--------------------------------------------------------------------
\section{Method}
%%--------------------------------------------------------------------

\subsection{Multi-Scale Shapelet Encoder}

We adopt the multi-scale shapelet representation of TimeCSL/CSL. For input
length $T$, $R$ scales are sampled uniformly in $[0.1T,0.8T]$, each with a fixed
number of shapelets; the encoder supports Euclidean distance, cosine similarity,
cross-correlation, and mixed measures. Let
$\mathcal{S}^{(r)}=\{s^{(r)}_{j}\}_{j=1}^{M_r}$ denote the shapelets at scale $r$.
The embedding concatenates per-scale responses,
\begin{equation}
z = \bigl[z^{(1)},z^{(2)},\dots,z^{(R)}\bigr],
\end{equation}
where $z^{(r)}$ collects the best-match responses to all shapelets at scale $r$.
Multi-scale shapelets capture local transients, meso-scale dynamics, and
long-range dependencies while remaining interpretable through explicit
subsequence patterns.

\subsection{Local Instance-Level Contrastive Learning}
\label{sec:local-cl}

Each client constructs two augmented views (noise, cropping, pooling,
quantization, time warping) following TS2Vec/TS-TCC~\cite{ts2vec,tstcc}. For
input $x$ with views $x_q,x_k$,
\begin{equation}
q = f(x_q;w_k),\qquad k = f(x_k;w_k).
\end{equation}
After normalization, the local InfoNCE-style objective is
\begin{equation}
\mathcal{L}_{\text{local}} = \mathrm{CE}\!\left(\frac{qk^\top}{\tau}, y\right),
\end{equation}
with temperature $\tau$ and positive-pair labels $y$; it encourages view
consistency and inter-sample separability.

\subsection{Local--Global Joint Contrastive Learning and Distillation}
\label{sec:joint-clkd}

Under non-IID FL, local instance contrast alone may optimize clients in
incompatible semantic coordinates~\cite{fedu2}. Inspired by MOON~\cite{moon},
FedX~\cite{fedx}, and FLESD~\cite{flesd}, we treat the previous global model as a
semantic teacher with representations
\begin{equation}
q_g = f(x_q; w_g),\qquad k_g = f(x_k; w_g),
\end{equation}
and add
\begin{equation}
\mathcal{L}_{\text{joint-CL}} = \mathrm{CE}\!\left(\frac{q_gk^\top}{\tau}, y\right),
\end{equation}
\begin{equation}
\mathcal{L}_{\text{joint-KD}} = D_{\mathrm{KL}}(q \,\|\, q_g) + D_{\mathrm{KL}}(k \,\|\, k_g),
\end{equation}
where $D_{\mathrm{KL}}$ is computed on similarity matrices. These terms preserve
local temporal adaptability while anchoring clients to a shared global semantic
space.

\subsection{Multi-Scale Structural Constraints and Scale-Level Distillation}
\label{sec:scale-constr}

Global-level distillation does not guarantee that each scale subspace is well
learned, so we align per scale:
\begin{equation}
\mathcal{L}_{\text{scale-CL}} = \sum_{r=1}^{R}\alpha_r \,
\mathrm{CE}\!\left(\frac{q_g^{(r)}(k^{(r)})^\top}{\tau}, y\right),
\label{eq:scale-cl}
\end{equation}
\begin{equation}
\mathcal{L}_{\text{scale-KD}} = \sum_{r=1}^{R}\alpha_r \left[
D_{\mathrm{KL}}\!\left(q^{(r)} \,\|\, q_g^{(r)}\right)+
D_{\mathrm{KL}}\!\left(k^{(r)} \,\|\, k_g^{(r)}\right)\right].
\end{equation}
The weight $\alpha_r$ is period-aware when \texttt{UseACF} is enabled
(Sec.~\ref{sec:period-score}) and uniform otherwise. Following CSL/TimeCSL,
covariance off-diagonal penalties on $Z^{(r)}$ reduce channel redundancy within
each scale, forming a structural alignment loss
$\mathcal{L}_{\text{structure}}$ that counteracts collapse and scale redundancy.

\subsection{Period-Aware Multi-Scale Weighting}
\label{sec:period-score}

Dominant periods and frequencies vary across datasets and channels, whereas the
encoder extracts patterns at $R$ discrete lengths. Equal weighting wastes
gradients on scales mismatched to local spectra. At each training step we call
\texttt{period\_score} on the current batch to produce non-negative weights
aligned with the $R$ reference scales, used as $\alpha_r$ in
Eq.~\eqref{eq:scale-cl} and in scale-level KD.

\paragraph{Reference scales and discrete anchors.}
Normalized anchors $\rho_r=\rho_{\min}+\frac{r-1}{R-1}(\rho_{\max}-\rho_{\min})$
with $\rho_{\min}=0.1$, $\rho_{\max}=0.8$ map to target lengths
$\tau_r=\lfloor \rho_r L \rfloor$, matching the multi-scale shapelet
configuration.

\paragraph{STFT branch (short--medium scales).}
For up to $k=\min(9,K)$ channels per sample we compute the STFT with $f_s=1$,
window $W=\max(4,\lfloor L/2\rfloor)$, $N_{\mathrm{fft}}=8W$, and hop
$h=\lfloor W/4\rfloor$. With power $P_c(f,t)=|Z_{xx}(f,t)|^2$, a 95th-percentile
mask selects high-energy bins and yields the masked mean frequency
\begin{equation}
\bar{f}_c=\frac{\sum_{f,t} f\,P_c(f,t)\,\mathbb{I}[\cdot]}
{\sum_{f,t} P_c(f,t)\,\mathbb{I}[\cdot]+\varepsilon}.
\end{equation}
The period proxy $s_c^{\mathrm{ST}}=1/(\bar{f}_c+\varepsilon)$ is soft-aligned to
the first $\lfloor R/2\rfloor$ anchors by Gaussian kernels with $\sigma=0.1L$,
accumulating into $\mathbf{S}^{\mathrm{ST}}$ on
$\mathcal{I}_{\mathrm{ST}}=\{1,\dots,\lfloor R/2\rfloor\}$.

\paragraph{ACF branch (medium--long scales).}
The sample autocorrelation $R_c(\ell)$ is computed with \texttt{fft=True} up to
lag $L-1$. Peaks in the second half, $\ell\ge\lfloor L/2\rfloor$ with
$R_c(\ell)>\theta_{\mathrm{acf}}=0.2$, vote inversely by distance to each anchor,
\begin{equation}
v_{c,r}(\ell)=\frac{1}{|\tau_r-\ell|+\varepsilon},\qquad
\mathbf{v}_c=\sum_{\ell\in\mathcal{P}_c}\bigl(v_{c,1}(\ell),\ldots,v_{c,R}(\ell)\bigr),
\end{equation}
yielding $\mathbf{S}^{\mathrm{ACF}}$ on
$\mathcal{I}_{\mathrm{ACF}}=\{\lfloor R/2\rfloor+1,\dots,R\}$.

\paragraph{Dual-branch fusion.}
The normalized STFT vector and the ACF subvector are fused with coefficient
$\beta$,
\begin{equation}
\boldsymbol{\pi}=\mathrm{Norm}\!\left(
\beta\,\mathrm{Embed}_{\mathcal{I}_{\mathrm{ST}}}\!\bigl(\hat{\mathbf{S}}^{\mathrm{ST}}\bigr)
+(1-\beta)\,\mathrm{Embed}_{\mathcal{I}_{\mathrm{ACF}}}\!\bigl(\hat{\mathbf{S}}^{\mathrm{ACF}}_{\mathcal{I}_{\mathrm{ACF}}}\bigr)
\right),
\label{eq:period-fuse}
\end{equation}
where $\mathrm{Embed}_{\mathcal{I}}$ pads a subvector back to $R$ dimensions and
$\mathrm{Norm}$ yields a non-negative simplex vector (uniform fallback if both
branches vanish). Hence each batch's spectral and ACF statistics set the
multi-scale loss emphasis without extra labels.

\subsection{Spilter: Adaptive Scale Stitching and Top-$m$ Selection}
\label{sec:spilter}

The multi-scale encoder decomposes into $R$ scale submodels with parameters
$\theta^{(r)}$, length $\ell_r$, and $M_r$ shapelets; the full model is
$w=\{\theta^{(1)},\dots,\theta^{(R)},\theta^{(\mathrm{shared})}\}$. Unlike
server-planned heterogeneous submodels, Spilter lets each client autonomously
select scales valuable under its local periodic/spectral structure and stitch
their features into a local submodel. Client $k$ activates
$\mathcal{R}_k\subseteq\{1,\dots,R\}$ with $|\mathcal{R}_k|=m<R$ and pulls,
trains, and uploads only $\{\theta^{(r)}:r\in\mathcal{R}_k\}$.

\paragraph{Local period-aware scoring.}
Client $k$ computes $\mathbf{s}_k=(s_{k,1},\dots,s_{k,R})$ with the STFT/ACF
mechanism of Sec.~\ref{sec:period-score}, reflecting the match between local
periodic/spectral structure and the anchors $\tau_r$. Scoring is fully
local---no raw data or system proxy is uploaded.

\paragraph{Distribution-aware top-$m$ selection.}
Given $\mathbf{s}_k$, client $k$ selects the top-$m$ scales
\begin{equation}
\mathcal{R}_k=\mathrm{TopM}(\mathbf{s}_k,m)
=\underset{\substack{\mathcal{R}\subseteq\{1,\dots,R\}\\|\mathcal{R}|=m}}{\arg\max}\;
\sum_{r\in\mathcal{R}} s_{k,r}.
\label{eq:spilter-topm}
\end{equation}
Scores are normalized and sorted stably, with ties preferring smaller scale
indices; the plan is fixed after one local precomputation. $m=R$ recovers
full-scale FedCSL and random scores yield the Random-top-$m$ ablation. Because
ranking uses only local distribution-aware scores, different clients naturally
obtain different $\mathcal{R}_k$, forming a multi-view ensemble.

\paragraph{Scale feature stitching.}
For input $x$, the selected features are concatenated,
\begin{equation}
z^{\mathrm{stitch}}_k = \bigl[z^{(r_1)} \,\|\, z^{(r_2)} \,\|\, \cdots \,\|\, z^{(r_m)}\bigr],
\qquad r_i\in\mathcal{R}_k .
\label{eq:stitch}
\end{equation}
Stitched representations drive $\mathcal{L}_{\text{local}}$,
$\mathcal{L}_{\text{joint-CL}}$, and $\mathcal{L}_{\text{joint-KD}}$; each active
scale additionally computes $\mathcal{L}_{\text{scale-CL}}$ and
$\mathcal{L}_{\text{scale-KD}}$ with period-aware weights. Clients thus train a
customized $m$-scale submodel: backbone losses on stitched features, auxiliary
alignment on component scales.

\paragraph{Coverage-aware scale allocation as a knapsack.}
Pure top-$m$ maximizes each client's local value but guarantees neither global
scale coverage nor local optimal total value. We therefore introduce binary
variables $x_{k,r}\in\{0,1\}$, a per-client memory budget (Sec.~\ref{sec:mast}),
and a minimum global coverage $N_{\min}$ per scale. Relaxing the coverage
equalities by a Lagrange multiplier $\mu_r\ge 0$ gives a decomposed objective,
\begin{equation}
\max_{\{x_{k,r}\}}\;\sum_{k=1}^{K}\sum_{r=1}^{R}
\bigl(s_{k,r}+\mu_r\bigr)x_{k,r}
\quad\text{s.t.}\quad
\sum_{r=1}^{R} c_{k,r}\,x_{k,r}\le B_k,\;\;
\sum_{r}x_{k,r}=m_k ,
\label{eq:knapsack}
\end{equation}
where $c_{k,r}$ is the estimated memory cost of scale $r$. Each client solves an
independent 0--1 knapsack with adjusted values $s_{k,r}+\mu_r$; the server
updates $\mu$ by subgradient steps, raising the value of under-covered scales
until coverage is restored. This yields memory-constrained, coverage-aware scale
trimming rather than a purely greedy top-$m$.

\subsection{Scale-Sharded Communication and Scale-Wise Aggregation}

In round $t$, client $k$ pulls
$\Theta_{k}^{\downarrow}(t)=\{\theta_g^{(r,t)}:r\in\mathcal{R}_k\}$ and uploads
$\Theta_{k}^{\uparrow}(t)=\{\theta_k^{(r,t+1)}:r\in\mathcal{R}_k\}$. The
bidirectional cost is
\begin{equation}
\mathcal{C}_{\mathrm{split}}(t)
=2\sum_{k=1}^{K}\sum_{r\in\mathcal{R}_k}B_r + O(KR),
\end{equation}
against the full-scale cost $\mathcal{C}_{\mathrm{full}}(t)=2K\sum_{r=1}^{R}B_r$;
with near-uniform per-scale payload,
$\mathcal{C}_{\mathrm{split}}/\mathcal{C}_{\mathrm{full}}\approx m/R$.
Computation scales analogously,
$\mathcal{F}_{\mathrm{split}}(t)=\sum_{k}E\,n_k\sum_{r\in\mathcal{R}_k}F_r$ versus
$\mathcal{F}_{\mathrm{full}}(t)=\sum_{k}E\,n_k\sum_{r}F_r$, so the saving depends
on the cost mix of the selected scales.

Let $\mathcal{K}_r(t)=\{k:r\in\mathcal{R}_k\}$. The server aggregates
\begin{equation}
\theta_g^{(r,t+1)}=
\begin{cases}
\displaystyle
\sum_{k\in\mathcal{K}_r(t)}
\frac{n_k}{\sum_{j\in\mathcal{K}_r(t)}n_j}\,
\theta_k^{(r,t+1)}, & \mathcal{K}_r(t)\neq\emptyset,\\[2ex]
\theta_g^{(r,t)}, & \mathcal{K}_r(t)=\emptyset,
\end{cases}
\label{eq:scale-agg}
\end{equation}
so untrained scales remain unchanged, consistent with heterogeneous FL such as
HeteroFL and FlexFL~\cite{heterofl,flexfl}.

\subsection{Overall Optimization Objective}
\label{sec:overall-obj}

We now assemble the loss components of
Secs.~\ref{sec:local-cl}--\ref{sec:scale-constr} into a single training
objective. Given its selected scales $\mathcal{R}_k$
(Eq.~\eqref{eq:spilter-topm}), client $k$ minimizes
\begin{equation}
\begin{aligned}
\mathcal{L}_k(w_k;\mathcal{R}_k)
={}& \gamma\,\mathcal{L}_{\mathrm{loc}}
+\lambda_5\,\mathcal{L}_{\mathrm{str}}
+\mathbb{I}_{\mathrm{Fed}}\Bigl[
\lambda_1\,\mathcal{L}_{\mathrm{jCL}}
+\lambda_2\,\mathcal{L}_{\mathrm{jKD}} \\
&+\lambda_3\,\mathcal{L}_{\mathrm{sCL}}
+\lambda_4\,\mathcal{L}_{\mathrm{sKD}}
\Bigr],
\end{aligned}
\label{eq:client-obj}
\end{equation}
where $\mathcal{L}_{\mathrm{loc}}$, $\mathcal{L}_{\mathrm{str}}$,
$\mathcal{L}_{\mathrm{jCL}}$, $\mathcal{L}_{\mathrm{jKD}}$,
$\mathcal{L}_{\mathrm{sCL}}$, and $\mathcal{L}_{\mathrm{sKD}}$ denote the local
instance-level contrastive, structural alignment, joint contrastive, joint
distillation, scale-level contrastive, and scale-level distillation losses
defined in Secs.~\ref{sec:local-cl}--\ref{sec:scale-constr}. The backbone terms
$\mathcal{L}_{\mathrm{loc}}$, $\mathcal{L}_{\mathrm{jCL}}$, and
$\mathcal{L}_{\mathrm{jKD}}$ act on the stitched features
(Eq.~\eqref{eq:stitch}), whereas $\mathcal{L}_{\mathrm{sCL}}$ and
$\mathcal{L}_{\mathrm{sKD}}$ act on the $m$ active component scales and carry
the period-aware weights $\alpha_r$ (Sec.~\ref{sec:period-score}). A single
balance parameter $\gamma\in[0,1]$, with $\zeta=1-\gamma$, controls the
trade-off between local adaptation and global alignment: the joint terms are
scaled by $\gamma$ and the scale-level terms by $\zeta$, up to fixed positive
constants, and the indicator $\mathbb{I}_{\mathrm{Fed}}\in\{0,1\}$ gates the
FedCSL-specific terms ($0$ for the local-only ablation, which disables
\texttt{UseJointCL}, \texttt{UseJointKD}, \texttt{UseScaleCL}, and
\texttt{UseScaleKD}).

The server approximately minimizes the sample-weighted sum of the client
objectives,
\begin{equation}
\mathcal{L}_g(w_g;\mathcal{R})
=\sum_{k=1}^{K}\frac{n_k}{N}\,\mathcal{L}_k(w_g;\mathcal{R}_k),
\label{eq:server-obj}
\end{equation}
which separates across scales under the aggregation rule of
Eq.~\eqref{eq:scale-agg}: each $\theta_g^{(r)}$ is updated as the
sample-weighted average over $\mathcal{K}_r(t)$, so the server step reduces to
$R$ independent per-scale subproblems solved in closed form each round.

The overall training problem is therefore a bilevel program over the scale
allocation $\mathcal{R}=\{\mathcal{R}_1,\dots,\mathcal{R}_K\}$,
\begin{equation}
\min_{\mathcal{R}}\;\sum_{k=1}^{K}\frac{n_k}{N}\,
\mathcal{L}_k\!\bigl(w_k^{\star}(\mathcal{R}_k);\mathcal{R}_k\bigr)
\quad\text{s.t.}\quad
|\mathcal{R}_k|=m,\;\;k=1,\dots,K,
\label{eq:overall-obj}
\end{equation}
where $w_k^{\star}(\mathcal{R}_k)=\arg\min_{w_k}\mathcal{L}_k(w_k;\mathcal{R}_k)$
is the local optimum of Eq.~\eqref{eq:client-obj}. The default allocator is the
distribution-aware top-$m$ rule of Eq.~\eqref{eq:spilter-topm}; when per-device
memory budgets $B_k$ and a global coverage floor $N_{\min}$ are enforced, the
allocation step is replaced by the Lagrangian-decomposed knapsack of
Eq.~\eqref{eq:knapsack}. Algorithm~\ref{alg:fedcsl} summarizes one training
round.

\begin{algorithm}[t]
\caption{FedCSL-Spilter training procedure}
\label{alg:fedcsl}
\begin{algorithmic}[1]
\State Initialize global model $w_g^{(0)}$
\For{each client $k$ \textbf{in parallel}}
    \State Compute period-aware scores $\mathbf{s}_k$ on local data
    \State Determine $\mathcal{R}_k\subseteq\{1,\dots,R\}$, $|\mathcal{R}_k|=m$
           via Eq.~\eqref{eq:spilter-topm}
\EndFor
\For{$t=0,1,\dots,T-1$}
    \For{each client $k$ \textbf{in parallel}}
        \State Pull $\Theta_k^{\downarrow}(t)=\{\theta_g^{(r,t)}:r\in\mathcal{R}_k\}$
        \State Augment inputs and stitch $|\mathcal{R}_k|$ scale features
               (Eq.~\eqref{eq:stitch})
        \State Compute local CL on stitched features; joint/scale CL/KD on
               $\mathcal{R}_k$
        \State Update local scale parameters; upload $\Theta_k^{\uparrow}(t)$
    \EndFor
    \State Server aggregates each scale via Eq.~\eqref{eq:scale-agg}
\EndFor
\end{algorithmic}
\end{algorithm}

%%--------------------------------------------------------------------
\section{Memory-Aware Multi-Scale Training with MAST}
\label{sec:mast}
%%--------------------------------------------------------------------

FedCSL-Spilter already contains two levels of memory optimization: the server
distills a useful scale set and distributes only those submodels, and each client
trims scales further under its own budget. A third level targets the local
multi-scale computation itself.

\subsection{The Activation Bottleneck}

Training proceeds in mini-batches, each with a forward and a backward pass. In
the forward pass every layer's output (activation) is kept for the backward
pass; once a layer's gradient is computed its activations are released. Peak GPU
memory therefore has three parts: model parameters, optimizer states, and
intermediate activations. Activations dominate. For a multivariate sample with
$D$ variables, sequence length $T$, and batch size $B$, a single scale submodule
with shapelet length $\ell$ and representation dimension $d$ stores
$O(B\,D\,T\,M_r\,d)$ values, and $R\times M_r$ submodules reside in memory
simultaneously---easily causing out-of-memory (OOM) errors. Unlike parameters and
optimizer states, activation size grows linearly with batch size, channel count,
and sequence length, which is precisely why large-scale MTS training is
memory-bound on 16--32\,GB GPUs.

\subsection{Checkpointing and Why Planning Is Hard}

Activation checkpointing keeps only selected activations in the forward pass and
recomputes the rest during backpropagation~\cite{chen2016checkpoint}. It is
supported natively (\texttt{torch.utils.checkpoint}), but default
implementations checkpoint whole layers indiscriminately. A good plan depends on
layer dependencies: if the activations of both $l_2$ and $l_3$ are dropped,
recomputing $l_3$ first requires recomputing $l_2$, producing a recursive
recomputation chain. Consequently, heuristic planners~\cite{mimose} lack global
optimality, while MILP planners such as
Checkmate~\cite{checkmate} and MONet~\cite{monet} must encode nonlinear
dependency constraints and may be slow to solve for very large networks.

\subsection{MAST: Exact Planning via Structural Decomposition}

MAST exploits two structural properties of the shapelet transformer encoder to
shrink the planning space to an exactly solvable one. First, memory is dominated
by submodules: the activations of a submodule $\phi_{r,m}$ with scale $r$ and
measure $m$ account for roughly $O(B\,D\,T\,M_r d)$, whereas the subsequent loss
computation needs only $O(B\,d\mathcal{R})$; for large $B$, $D$, and $T$ the gap
is decisive. Second, submodules are mutually independent: different
$\phi_{r,m}$ are only concatenated at the output and share no data dependency, so
any subset of activations can be dropped and recomputed without triggering
recursion.

Let $x_{r,m}\in\{0,1\}$ indicate whether submodule $\phi_{r,m}$ is checkpointed.
Writing $C_{r,m}$ for its forward cost (obtained by profiling once), the total
cost is
\begin{equation}
J=\sum_{r,m}\bigl[\,C_{r,m}\,x_{r,m} + C_{r,m}\,(1-x_{r,m})\cdot \rho_{r,m}\,\bigr],
\end{equation}
where $\rho_{r,m}$ is the recomputation factor, so the optimization reduces to
minimizing the recomputation overhead
\begin{equation}
\min_{\{x_{r,m}\}}\;\sum_{r,m} C_{r,m}\,(1-x_{r,m})\,\rho_{r,m}
\end{equation}
subject to the peak-memory constraints of the five phases---forward submodule,
forward loss, backward loss, backward submodule, and submodule recomputation---
each staying below the budget. All symbols except the decisions are constants,
so this is a \emph{binary linear program} solvable exactly by Gurobi.

\subsection{Using MAST Inside Federated Rounds}

MAST does not change the mathematical model and therefore does not affect
accuracy: it is a memory--time regulator. In synchronous FL the round duration is
set by the straggler, so a faster client can spend its spare memory to train with
fewer checkpoints (more activation memory, less recomputation) without altering
what is learned. Conversely, a memory-constrained client runs a more aggressive
plan and remains in the same round. The per-client budgets $B_k$ used in the
scale knapsack (Eq.~\eqref{eq:knapsack}) come directly from this planner.

%%--------------------------------------------------------------------
\section{Experimental Setup}
%%--------------------------------------------------------------------

\subsection{Datasets and Federated Partitioning}

We use five MTS classification benchmarks:
\begin{itemize}
  \item \textbf{HAR} (UCI Human Activity Recognition): 6-class smartphone
        inertial sensing, length 128, 9 channels; used for multi-$\alpha$
        tables, convergence curves, system efficiency, and ablations.
  \item \textbf{FaceDetection} (UEA): length 62, 3 channels, 2 classes.
  \item \textbf{LSST} (UEA): astronomical light curves, length 36, 6 channels,
        14 classes.
  \item \textbf{Epilepsy} (UEA): length 206, 3 channels, 4 classes.
  \item \textbf{SleepEDF}: multi-channel EEG/EOG sleep staging, 5 stages.
\end{itemize}
FaceDetection, LSST, and HAR sweep the Dirichlet concentration
$\alpha\in\{0.05,0.1,0.3,0.5,1.0,10.0\}$; Epilepsy and SleepEDF ablations use the
default non-IID split. MAST is evaluated on four large UEA datasets: Cricket
(Cr), DuckDuckGeese (DD), MotorImagery (MI), and PEMS-SF (PE).

\subsection{Baselines}

\textbf{Tier 1: federated representation baselines and our methods.} All methods
share the multi-scale shapelet encoder and the SVM evaluation, varying only the
FL objective and aggregation: (i) classic FL---FedAvg~\cite{fedavg} and
FedProx~\cite{fedprox}; (ii) federated self-supervised baselines---FedBYOL,
FedU2~\cite{fedu2}, and Orchestra; (iii) full-scale FedCSL with joint/scale CL/KD
and STFT/ACF weighting over all $R$ scales; (iv) FedCSL-Spilter with local-score
top-$m$ stitching, $m\in\{1,2,4\}$.

\textbf{Tier 2: module ablations.} Local-only contrastive learning, FedAvg-style
aggregation (AVG), Federated Contrastive All (full FedCSL modules), and
Federated Contrastive + multi-scale adaptation (Spilter), plus Random-top-$m$
variants that replace period scores with uniform random scores.

\subsection{Training Configuration}

Each round, clients determine $\mathcal{R}_k$ via Eq.~\eqref{eq:spilter-topm},
pull the corresponding scale parameters, train locally for several epochs, and
upload the trained scales; the server aggregates via Eq.~\eqref{eq:scale-agg}. We
use SGD, mixed distance measures, and unified hyperparameters in configuration
files; Spilter runs set \texttt{allocation\_mode=local\_score\_topm}. Every
algorithm has its own YAML configuration and launch script, and the batch
$\alpha$-matrix experiments are dispatched across multiple GPUs.

\subsection{Evaluation Protocol}

Following TimeCSL~\cite{timescsl} and the linear evaluation of
TS2Vec/TS-TCC~\cite{ts2vec,tstcc}, we extract train/test representations with the
global shapelet encoder and train an SVM classifier, decoupling representation
quality from end-to-end classification. We report the mean test SVM accuracy over
the last 10 rounds, together with average rank, rounds-to-convergence, and
wall-clock and memory statistics.

%%--------------------------------------------------------------------
\section{Results and Discussion}
%%--------------------------------------------------------------------

\subsection{Multi-Scale Comparative Learning}

Table~\ref{tab:face_dirichlet} reports FaceDetection accuracy under six Dirichlet
concentrations. FedCSL attains the best or second-best result in all six
settings, staying within $0.0144$ of its own range ($0.5626$--$0.5770$), whereas
FedAvg spans $0.5376$--$0.5633$ and degrades clearly under strong heterogeneity.
The gap between $\alpha=0.05$ and $\alpha=10.0$ is only $0.0084$ for FedCSL,
whereas FedAvg and FedProx both drop when $\alpha$ falls from $0.1$ to $0.05$,
confirming their reliance on the IID assumption. This robustness follows from the
multi-scale shapelet encoder: submodules at different temporal scales capture
patterns from local to global, so features remain stable when client
distributions shift. FedCSL peaks at $0.5770$ for $\alpha=1.0$ and has its
largest advantage over FedAvg at $\alpha=0.5$ ($+0.0351$, about 6.5\%), which
suggests that moderate heterogeneity supplies enough negative-sample diversity
for contrastive learning without destroying semantic consistency. Among
unsupervised baselines, Orchestra and FedU2 are stable but remain $0.01$--$0.02$
below FedCSL; FedBYOL reaches $0.5763$ at $\alpha=1.0$ but falls to $0.5530$
under strong heterogeneity ($\alpha=0.05$). The joint objective of multi-grained
contrast plus multi-scale alignment is therefore better suited to federated
heterogeneity than a single global contrastive loss.

\begin{table*}[t]
  \centering
  \caption{FaceDetection accuracy under varying Dirichlet non-IID strength
  (mean test SVM accuracy over the last 10 rounds; bold = best per column).}
  \label{tab:face_dirichlet}
  \begin{tabular}{lcccccc}
    \toprule
    Method & $\alpha{=}0.05$ & $\alpha{=}0.1$ & $\alpha{=}0.3$ & $\alpha{=}0.5$ & $\alpha{=}1.0$ & $\alpha{=}10.0$ \\
    \midrule
    FedAvg   & 0.5594 & 0.5633 & 0.5459 & 0.5386 & 0.5501 & 0.5376 \\
    FedProx  & 0.5430 & 0.5438 & 0.5486 & 0.5508 &        & 0.5336 \\
    FedBYOL  & 0.5530 & 0.5592 & 0.5627 & 0.5616 & 0.5763 & 0.5578 \\
    FedU2    & 0.5621 & 0.5634 & 0.5463 & 0.5630 & 0.5550 & 0.5618 \\
    Orchestra& 0.5609 & 0.5530 & 0.5646 & 0.5620 & 0.5591 & \textbf{0.5632} \\
    FedCSL   & \textbf{0.5710} & \textbf{0.5646} & \textbf{0.5701} & \textbf{0.5737} & \textbf{0.5770} & 0.5626 \\
    \bottomrule
  \end{tabular}
\end{table*}

\subsection{Non-IID Comparison on HAR}

Table~\ref{tab:har_dirichlet} reports nine methods over six values of $\alpha$ on
HAR. Spilter-$m{=}2$ and Spilter-$m{=}4$ outperform full-scale FedCSL and the
self-supervised baselines at most concentrations: Spilter-$m{=}2$ is best at
$\alpha=0.1,1.0,10.0$ (reaching $0.9655$ at $\alpha=10.0$), while Spilter-$m{=}4$
is best at $\alpha=0.3$ and $0.5$ ($0.9521$ and $0.9586$) and at $\alpha=0.1$
($0.9446$). Spilter-$m{=}1$ is slightly less accurate but has the lowest memory,
forming an explicit accuracy--cost trade-off. Notably, stitched submodels beating
full-scale FedCSL under moderate to strong non-IID data indicates a structured
regularization effect: restricting each client to a subset of scales prevents
gradients from being diluted across scale subspaces that the local data cannot
support.

\begin{table*}[t]
  \centering
  \caption{HAR accuracy under varying Dirichlet non-IID strength (mean test SVM
  accuracy over the last 10 rounds; bold = best per column).}
  \label{tab:har_dirichlet}
  \begin{tabular}{lcccccc}
    \toprule
    Method & $\alpha{=}0.05$ & $\alpha{=}0.1$ & $\alpha{=}0.3$ & $\alpha{=}0.5$ & $\alpha{=}1.0$ & $\alpha{=}10.0$ \\
    \midrule
    FedAvg            & 0.9332 & 0.9370 & 0.9338 & 0.9462 & 0.9445 & 0.9484 \\
    FedProx           & \textbf{0.9472} & 0.9399 & 0.9346 & 0.9384 & 0.9398 & 0.9458 \\
    FedBYOL           & 0.9268 & 0.9285 & 0.9336 & 0.9284 & 0.9331 & 0.9267 \\
    FedU2             & 0.9252 & 0.9240 & 0.9307 & 0.9286 & 0.9317 & 0.9327 \\
    Orchestra         & 0.9344 & 0.9336 & 0.9293 & 0.9315 & 0.9325 & 0.9333 \\
    FedCSL            & 0.9315 & 0.9339 & 0.9306 & 0.9382 & 0.9475 & 0.9514 \\
    \midrule
    Spilter ($m{=}1$) & 0.9261 & 0.9236 & 0.9350 & 0.9378 & 0.9253 & 0.9331 \\
    Spilter ($m{=}2$) & 0.9377 & 0.9346 & 0.9489 & 0.9527 & 0.9405 & \textbf{0.9655} \\
    Spilter ($m{=}4$) & 0.9414 & \textbf{0.9446} & \textbf{0.9521} & \textbf{0.9586} & \textbf{0.9573} & 0.9630 \\
    \bottomrule
  \end{tabular}
\end{table*}

\subsection{LSST and Ablation Study}

Table~\ref{tab:lsst_dirichlet} reports LSST, where astronomical light curves are
short (length 36) and highly heterogeneous. FedU2 is strongest on this benchmark
and Spilter-$m{=}4$ approaches it ($0.6081$ at $\alpha=10.0$) but does not
surpass it---an honest limitation we discuss in Sec.~\ref{sec:limitations}.
Table~\ref{tab:ablation} isolates the contribution of each module. Moving from
FedAvg-style aggregation (AVG) to the full FedCSL objective, and finally to
Spilter on top of it, yields consistent gains on Epilepsy, HAR, and SleepEDF; the
largest improvement appears on SleepEDF
($0.7797\rightarrow0.8052\rightarrow0.8213$), where the inter-client stage
distribution is most skewed. The local-only column is not reported in this run
and is shown as a placeholder.

\begin{table*}[t]
  \centering
  \caption{LSST accuracy under varying Dirichlet non-IID strength (mean test SVM
  accuracy over the last 10 rounds; bold = best per column).}
  \label{tab:lsst_dirichlet}
  \begin{tabular}{lcccccc}
    \toprule
    Method & $\alpha{=}0.05$ & $\alpha{=}0.1$ & $\alpha{=}0.3$ & $\alpha{=}0.5$ & $\alpha{=}1.0$ & $\alpha{=}10.0$ \\
    \midrule
    FedAvg            & 0.5857 & 0.5833 & 0.5901 & 0.5716 & 0.5770 & 0.5739 \\
    FedProx           & 0.5867 & 0.5951 & 0.5793 & 0.5777 & 0.5781 & 0.5794 \\
    FedBYOL           & 0.6033 & 0.6092 & 0.6071 & 0.6056 & 0.6051 & 0.6025 \\
    FedU2             & \textbf{0.6103} & \textbf{0.6102} & \textbf{0.6115} & \textbf{0.6185} & \textbf{0.6201} & \textbf{0.6138} \\
    Orchestra         & 0.6068 & 0.6066 & 0.6084 & 0.6095 & 0.6020 & 0.6018 \\
    FedCSL            & 0.5978 & 0.6054 & 0.5992 & 0.5926 & 0.6008 & 0.5968 \\
    \midrule
    Spilter ($m{=}4$) & 0.6013 & 0.5995 & 0.5897 & 0.6052 & 0.5940 & 0.6081 \\
    \bottomrule
  \end{tabular}
\end{table*}

\begin{table}[t]
  \centering
  \caption{Module ablation on Epilepsy, HAR, and SleepEDF (mean test SVM
  accuracy).}
  \label{tab:ablation}
  \resizebox{\columnwidth}{!}{%
  \begin{tabular}{lcccc}
    \toprule
    Dataset & Local & AVG & Fed.\ Contrast.\ All & Fed.\ Contrast. \\
            &       &     &                      & + Multi-Scale Adapt. \\
    \midrule
    Epilepsy & --- & 0.9804 & 0.9839 & 0.9852 \\
    HAR      & --- & 0.9281 & 0.9342 & 0.9338 \\
    SleepEDF & --- & 0.7797 & 0.8052 & 0.8213 \\
    \bottomrule
  \end{tabular}%
  }
\end{table}

\subsection{System Efficiency and Client Memory}
\label{sec:efficiency}

Table~\ref{tab:har_memory} reports HAR micro-benchmarks ($\alpha=0.1$, $K=10$,
batch size 32, one local epoch, RTX 3090), measuring peak GPU memory with
\texttt{torch.cuda.max\_memory\_allocated()} after warmup. FedCSL activates all
$R=8$ scales; Spilter trains stitched $m$-scale submodels and reduces memory
almost linearly in $m$. Scale stitching alone cuts mean peak memory from
$178.1$\,MB to $51.5$\,MB at $m{=}1$ (71.1\% saving, $3.46\times$ compression),
while Spilter-$m{=}4$ stays Pareto-competitive on accuracy
(Table~\ref{tab:har_dirichlet}). The slowest epoch also drops from
$4.70$\,s to $3.17$\,s ($1.48\times$ speedup) at $m{=}1$, and Spilter-$m{=}4$
reaches $0.9446$ accuracy---above FedCSL's $0.9339$---at $3.76$\,s. Over 200
rounds the local-compute lower bound falls to about 68\% of FedCSL for $m{=}1$,
which is the quantity that matters when the round makespan is determined by the
straggler.

\begin{table*}[t]
  \centering
  \caption{HAR client peak memory and slowest-client epoch time ($\alpha=0.1$,
  RTX 3090, $K=10$). Mem Saving $=1-\bar{M}/\bar{M}_{\mathrm{FedCSL}}$;
  Compress. Ratio $=\bar{M}_{\mathrm{FedCSL}}/\bar{M}$. Bold: Spilter savings
  versus FedCSL.}
  \label{tab:har_memory}
  \begin{tabular}{lcccccc}
    \toprule
    Method & $m$ & Mean Peak (MB) $\downarrow$ & Max Peak (MB) $\downarrow$
           & Mem Saving $\uparrow$ & Compress. Ratio $\uparrow$ & Slowest Epoch (s) $\downarrow$ \\
    \midrule
    FedAvg            & 8 & 177.0 & 180.3 & 0.6\%  & 1.01$\times$ & \textbf{1.28} \\
    FedProx           & 8 & 179.2 & 182.4 & $-$0.6\% & 0.99$\times$ & 1.47 \\
    FedBYOL           & 8 & 182.9 & 183.4 & $-$2.7\% & 0.97$\times$ & 1.89 \\
    FedU2             & 8 & 182.9 & 183.4 & $-$2.7\% & 0.97$\times$ & 1.93 \\
    Orchestra         & 8 & 231.3 & 233.5 & $-$29.9\% & 0.77$\times$ & 2.28 \\
    FedCSL (baseline) & 8 & 178.1 & 181.3 & --- & 1.00$\times$ (ref) & 4.70 \\
    \midrule
    Spilter-$m{=}1$   & 1 & \textbf{51.5} & \textbf{51.7} & \textbf{71.1\%} & \textbf{3.46$\times$} & 3.17 \\
    Spilter-$m{=}2$   & 2 & \textbf{78.7} & \textbf{79.1} & \textbf{55.8\%} & \textbf{2.26$\times$} & 3.43 \\
    Spilter-$m{=}4$   & 4 & \textbf{130.3} & \textbf{131.5} & \textbf{26.8\%} & \textbf{1.37$\times$} & 3.76 \\
    \bottomrule
  \end{tabular}
\end{table*}

\begin{figure}[t]
  \centering
  \includegraphics[width=0.98\linewidth]{figures/convergence.png}
  \caption{HAR test accuracy versus communication rounds for varying Dirichlet
  $\alpha$ ($3\times2$ panels, $\alpha=0.05,0.1,0.3,0.5,1.0,10.0$ left-to-right,
  top-to-bottom). Curves show Spilter with \texttt{local\_top\_m} $\in\{1,2,4\}$
  against FedCSL; y-axes are zoomed per panel.}
  \label{fig:dirichlet}
\end{figure}

\subsection{Resource-Imbalanced Deployment}
\label{sec:resource-hetero}

In real deployments client memory and compute are uneven. Sec.~\ref{sec:spilter}
keeps scale ranking based on purely local period scores; device heterogeneity is
then modeled by \emph{lowering each client's scale budget}: strong devices use
$m{=}4$, medium devices $m{=}2$, weak devices $m{=}1$, giving a cardinality
degradation chain $m=4\rightarrow2\rightarrow1$. This matches
Table~\ref{tab:har_memory}: mean peak memory and per-epoch time decrease roughly
monotonically with $m$ (mean peak memory $130.3$, $78.7$, and $51.5$\,MB
for $m=4,2,1$, and mean epoch time $1.62$, $1.44$, and $1.34$\,s
epoch), while Table~\ref{tab:har_dirichlet} shows accuracy degrades slowly; at
$\alpha{=}0.1$ the three levels reach $0.9446$, $0.9346$, and $0.9236$. In other
words, dropping from $m{=}4$ to $m{=}1$ reduces mean peak memory to about 40\%
and mean epoch time by roughly 17\%, at an accuracy cost of about 2.1 points.

\paragraph{Local-score top-$m$ versus Random-top-$m$.}
To isolate whether period-aware scoring beats arbitrary scale subsets, we compare
local-score top-$m$ (Eq.~\eqref{eq:spilter-topm}) with Random-top-$m$ (period
scores replaced by uniform random scores) under the same $m{=}4$ and the same
federated splits. Table~\ref{tab:random_vs_topm} shows that local-score top-$m$
wins by 1.3--4.8 percentage points at every $\alpha$, with the largest gaps under
strong non-IID ($\alpha=0.05$: $0.9372$ vs.\ $0.9238$; $\alpha=0.1$: $0.9411$
vs.\ $0.9247$). Local STFT/ACF scores therefore select scale combinations
aligned with each client's dominant periods instead of arbitrary subsets.

\begin{table}[t]
  \centering
  \caption{HAR accuracy of Local-score top-$m$ versus Random-top-$m$ ($m{=}4$;
  mean test SVM accuracy over the last 10 rounds; bold: better at the same
  $\alpha$).}
  \label{tab:random_vs_topm}
  \begin{tabular}{lcc}
    \toprule
    $\alpha$ & Random-top-$m$ & Local-score top-$m$ \\
    \midrule
    0.05  & 0.9238 & \textbf{0.9372} \\
    0.1   & 0.9247 & \textbf{0.9411} \\
    0.3   & 0.9460 & \textbf{0.9520} \\
    0.5   & 0.9482 & \textbf{0.9513} \\
    1.0   & 0.9497 & \textbf{0.9529} \\
    10.0  & 0.9529 & \textbf{0.9577} \\
    \bottomrule
  \end{tabular}
\end{table}

\paragraph{Coverage-aware knapsack.}
The experiments above use greedy top-$m$. The knapsack formulation
(Eq.~\eqref{eq:knapsack}) is the natural next step when per-device memory budgets
differ and some scales would otherwise be covered by too few clients. Because the
Lagrangian decomposition leaves each client with a standard 0--1 knapsack, the
per-round overhead is small, and the multiplier update lets the server steer
coverage without inspecting local data.

\subsection{MAST: Memory--Time Trade-off}

MAST is evaluated on four UEA datasets with large channel counts and sequence
lengths: Cricket, DuckDuckGeese, MotorImagery, and PEMS-SF. Across all planners
and budgets the normalized time stays below 1 (all are faster than full-model
PyTorch checkpointing) and actual memory consumption always stays under the
budget, confirming that planning under a hard memory cap is effective. MAST is
faster than every baseline, achieving $1.21\times$--$2.72\times$ speedup over
PyTorch. Against the state-of-the-art planners Mimose, Checkmate, and MONet, MAST
reduces training time by 1.43\%--7.65\% at the same budget, and saves
6.4\%--21.9\% of GPU memory at the same computational cost (estimated by a
quadratic fit of the time--budget curve).

To examine this more closely we profiled MotorImagery over time under a 16\,GB
budget using the plans produced by MAST and by Checkmate (the second-fastest
planner). In the first iteration the forward phases are close: both take about
0--350\,ms and Checkmate's peak is slightly lower. In the backward phase,
however, Checkmate releases less memory, particularly between 650 and 750\,ms,
so its global peak reaches 14.7\,GB, 4.7\% higher than MAST; its backward pass
also takes longer, indicating heavier recomputation. This validates the
customized formulation.

These gains matter most for large training data and iterative development, where
GPU time is expensive (for reference, a 16\,GB V100 costs \$2.48 per hour on
Google Cloud). Beyond speed, MAST provides the per-client budget $B_k$ used by
the scale knapsack, closing the loop between the systems layer and the
representation layer.

\subsection{Interpretability Analysis}

Because representations are built from explicit subsequence patterns, they
support visualization that black-box CNN/Transformer embeddings do not: class-wise
embedding structure, per-scale response strength, selected scale sets
$\mathcal{R}_k$ under different local periodicities, and the trajectory of the
period-aware weights $\boldsymbol{\pi}$ during training. In practice, inspecting
$\mathcal{R}_k$ shows which temporal scales a client's data actually supports,
which is useful both for debugging federated partitions and for explaining why a
given client benefits from stitching.

\subsection{Limitations and Ongoing Work}
\label{sec:limitations}

Several gaps remain, and we list them explicitly.
(i)~Communication and system experiments do not yet cover all five datasets; the
efficiency study concentrates on HAR.
(ii)~Mixed per-client $m_k$ within the same round has not been evaluated; the
current study uses a uniform degradation chain, and the coverage-aware knapsack
is formulated but not yet deployed at scale.
(iii)~Scale selection is precomputed and fixed before training; an online
top-$m$ strategy that updates $\mathcal{R}_k$ across rounds is not yet realized.
(iv)~FedCSL-Spilter does not surpass FedU2 on LSST, which needs further analysis
and hyperparameter tuning; short light curves may not expose enough multi-scale
periodic structure.
(v)~Evaluation focuses on classification; clustering and anomaly detection are
future work~\cite{ts2vec,timescsl}.

%%--------------------------------------------------------------------
\section{Conclusion}
%%--------------------------------------------------------------------

We presented FedCSL-Spilter for non-IID multivariate time series, coupling
TimeCSL/CSL multi-scale shapelet learning with client-autonomous scale stitching:
local STFT/ACF scores select the top-$m$ scales, shapelet features are stitched
into local submodels, clients upload scale-sharded updates, and the server
aggregates scale-wise. Local--global joint CL/KD, multi-scale structural
alignment, and period-aware weighting mitigate client drift; a
covering-constrained knapsack extends greedy top-$m$ to memory-budgeted,
coverage-aware allocation; and MAST turns activation-checkpoint planning into an
exact binary linear program that regulates the memory--time trade-off per device
without affecting accuracy. Experiments on HAR, FaceDetection, LSST, Epilepsy,
SleepEDF, and four large UEA datasets show favorable
accuracy--memory--latency trade-offs, with up to 71.1\% memory savings,
$1.48\times$ faster slowest-client epochs, and 1.43\%--7.65\% lower training time
than state-of-the-art checkpoint planners at equal budgets. Future work includes
mixed $m_k$ participation, online scale selection, and personalized federated
extensions.

\section*{Acknowledgments}

The author thanks the advisor and the research group for discussions on
federated learning and memory-efficient training. Funding and compute-resource
acknowledgments will be added before submission.

\bibliographystyle{ACM-Reference-Format}
\bibliography{refs,refs_acm}

\end{document}
